\documentclass[12pt,a4paper]{article}
\usepackage[margin=2.5cm]{geometry}
\usepackage[T1]{fontenc}
\usepackage[utf8]{inputenc}
\usepackage[spanish]{babel}
\usepackage{enumitem}
\usepackage{hyperref}
\usepackage{longtable}
\usepackage{array}
\usepackage{setspace}

\setlength{\parindent}{0pt}
\setlength{\parskip}{0.6em}
\setlist[itemize]{leftmargin=1.5em}
\setlist[enumerate]{leftmargin=1.8em}
\onehalfspacing

\begin{document}

\begin{center}
{\Large \textbf{Documentaci\'on t\'ecnica}}\\[0.5cm]
{\LARGE \textbf{Instructivo para el cliente}}\\[0.8cm]
{\large Martinez, Nahuel}\\[0.2cm]
{\normalsize \texttt{jnahuel.developer@gmail.com}}\\[0.2cm]
{\normalsize 27/03/2026}
\end{center}

\vspace{0.8cm}

\section*{1. Objetivo}

Este documento explica c\'omo funciona la aplicaci\'on \textbf{Gesti\'on de Fardos}, qu\'e opciones de configuraci\'on utiliza y cu\'ales son las funciones disponibles para la operaci\'on diaria.

La idea es que el usuario pueda:

\begin{itemize}
\item entender qu\'e hace el sistema;
\item configurar correctamente la balanza y el pulsador;
\item operar la pantalla principal;
\item acceder al modo t\'ecnico cuando sea necesario;
\item identificar d\'onde quedan guardados los registros y los archivos exportados.
\end{itemize}

\section*{2. L\'ogica general del programa}

La aplicaci\'on trabaja con dos dispositivos conectados por puertos serie:

\begin{itemize}
\item \textbf{Balanza}: env\'ia continuamente el peso.
\item \textbf{Pulsador}: al presionarse, env\'ia la trama \texttt{\$P1!}.
\end{itemize}

La l\'ogica de funcionamiento es la siguiente:

\begin{enumerate}
\item La aplicaci\'on lee el peso actual enviado por la balanza.
\item Cuando el pulsador env\'ia \texttt{\$P1!}, la aplicaci\'on responde \texttt{\$B1!}.
\item En ese momento toma el \'ultimo peso v\'alido disponible.
\item Si el peso est\'a dentro del rango permitido, se guarda un registro en la base local.
\item Si el peso est\'a fuera de rango o no hay una lectura v\'alida, la captura se rechaza y se informa en pantalla.
\end{enumerate}

De esta forma, la balanza informa el peso y el pulsador define el momento exacto en el que se debe registrar la pesada.

\section*{3. Funciones disponibles}

\subsection*{3.1 Pantalla principal}

La pantalla principal muestra:

\begin{itemize}
\item el peso actual;
\item el estado de la \'ultima opresi\'on del pulsador;
\item el \'ultimo registro guardado;
\item el acceso a edici\'on de registros;
\item el acceso a exportaci\'on a Excel.
\end{itemize}

\subsection*{3.2 Guardado autom\'atico}

Cada vez que el pulsador env\'ia \texttt{\$P1!}, la aplicaci\'on intenta guardar una pesada.

El guardado solo se realiza si:

\begin{itemize}
\item la balanza entreg\'o un valor interpretable;
\item el peso est\'a dentro del rango configurado.
\end{itemize}

\subsection*{3.3 Edici\'on de registros}

Desde la pantalla principal se puede editar una pesada ya guardada.

El flujo es:

\begin{enumerate}
\item presionar \textbf{Editar registro};
\item elegir el n\'umero de pesada;
\item ingresar la contrase\~na de edici\'on;
\item confirmar la operaci\'on.
\end{enumerate}

Al editar una pesada:

\begin{itemize}
\item el registro no se elimina;
\item el peso queda en \texttt{0};
\item la modificaci\'on queda registrada.
\end{itemize}

\subsection*{3.4 Exportaci\'on a Excel}

Desde la pantalla principal se puede exportar un rango de fechas a un archivo \texttt{.xlsx}.

La exportaci\'on incluye:

\begin{itemize}
\item n\'umero de fardo;
\item d\'ia;
\item hora;
\item kg.
\end{itemize}

Los archivos se guardan en la carpeta definida en la configuraci\'on.

\subsection*{3.5 Modo Service}

El modo Service permite realizar diagn\'ostico t\'ecnico.

Se abre con:

\begin{itemize}
\item \texttt{Ctrl + Shift + S}
\item \texttt{Ctrl + Alt + Shift + S}
\end{itemize}

Sus funciones principales son:

\begin{itemize}
\item ver el estado de conexi\'on de balanza y pulsador;
\item ver datos recibidos en crudo;
\item ver la interpretaci\'on del peso y la tara;
\item revisar errores de comunicaci\'on;
\item borrar registros hist\'oricos por fecha.
\end{itemize}

\section*{4. Archivo de configuraci\'on}

La aplicaci\'on utiliza un archivo llamado \texttt{config.json}, ubicado junto al ejecutable.

Si se modifica este archivo, es necesario cerrar y volver a abrir la aplicaci\'on.

\subsection*{4.1 Estructura general}

Las secciones principales son:

\begin{itemize}
\item \texttt{Scale}
\item \texttt{Button}
\item \texttt{Database}
\item \texttt{Thresholds}
\item \texttt{Passwords}
\item \texttt{Export}
\end{itemize}

\subsection*{4.2 Configuraci\'on de la balanza}

La secci\'on \texttt{Scale} define dos cosas:

\begin{itemize}
\item c\'omo abrir el puerto serie;
\item c\'omo interpretar el dato recibido.
\end{itemize}

Campos principales:

\begin{longtable}{|>{\raggedright\arraybackslash}p{4cm}|>{\raggedright\arraybackslash}p{10cm}|}
\hline
\textbf{Campo} & \textbf{Descripci\'on} \\
\hline
\texttt{Protocol} & Protocolo de datos de la balanza. Actualmente el valor recomendado es \texttt{w180-t}. \\
\hline
\texttt{PortName} & Puerto COM de la balanza. \\
\hline
\texttt{BaudRate} & Velocidad del puerto serie. \\
\hline
\texttt{DataBits} & Cantidad de bits de datos. \\
\hline
\texttt{Parity} & Paridad del puerto. \\
\hline
\texttt{StopBits} & Bits de stop. \\
\hline
\texttt{Handshake} & Control de flujo. \\
\hline
\texttt{NewLine} & Solo aplica si se usa un protocolo por l\'inea, como \texttt{simple-ascii}. \\
\hline
\texttt{WeightDecimalDigits} & Define cu\'antos d\'igitos del valor recibido corresponden a la parte decimal del peso. \\
\hline
\texttt{TareDecimalDigits} & Define cu\'antos d\'igitos del valor recibido corresponden a la parte decimal de la tara. \\
\hline
\end{longtable}

\subsection*{4.3 Importancia de \texttt{WeightDecimalDigits} y \texttt{TareDecimalDigits}}

La balanza entrega el valor como una secuencia de d\'igitos enteros, por ejemplo:

\begin{itemize}
\item \texttt{099500}
\item \texttt{100250}
\item \texttt{000000}
\end{itemize}

Como el equipo no informa d\'onde est\'a la coma decimal, la aplicaci\'on ahora lo toma desde configuraci\'on.

Ejemplos:

\begin{itemize}
\item si \texttt{WeightDecimalDigits = 3}, el valor \texttt{099500} se interpreta como \texttt{99,500 kg};
\item si \texttt{WeightDecimalDigits = 2}, el valor \texttt{099500} se interpreta como \texttt{995,00 kg}.
\end{itemize}

Lo mismo aplica para la tara mediante \texttt{TareDecimalDigits}.

\subsection*{4.4 Configuraci\'on del pulsador}

La secci\'on \texttt{Button} define el puerto del pulsador.

Campos principales:

\begin{itemize}
\item \texttt{PortName}
\item \texttt{BaudRate}
\item \texttt{DataBits}
\item \texttt{Parity}
\item \texttt{StopBits}
\item \texttt{Handshake}
\end{itemize}

El pulsador debe usar un puerto distinto al de la balanza.

\subsection*{4.5 Base local}

La secci\'on \texttt{Database} define la ubicaci\'on del archivo de base SQLite:

\begin{itemize}
\item \texttt{FilePath}
\end{itemize}

All\'i quedan guardadas las pesadas.

\subsection*{4.6 Rango de pesos}

La secci\'on \texttt{Thresholds} define el rango v\'alido de guardado:

\begin{itemize}
\item \texttt{MinKg}
\item \texttt{MaxKg}
\end{itemize}

Si una pesada queda fuera de ese rango, no se registra.

\subsection*{4.7 Contrase\~nas}

La secci\'on \texttt{Passwords} contiene:

\begin{itemize}
\item \texttt{Edit}: contrase\~na para editar registros;
\item \texttt{Service}: contrase\~na para ingresar al modo Service.
\end{itemize}

\subsection*{4.8 Carpeta de exportaci\'on}

La secci\'on \texttt{Export} define:

\begin{itemize}
\item \texttt{Folder}: carpeta donde se generan los archivos Excel.
\end{itemize}

\section*{5. Ejemplo de configuraci\'on}

\begin{verbatim}
"Scale": {
  "Protocol": "w180-t",
  "PortName": "COM4",
  "BaudRate": 9600,
  "DataBits": 7,
  "Parity": "Even",
  "StopBits": "Two",
  "Handshake": "None",
  "NewLine": "\n",
  "WeightDecimalDigits": 3,
  "TareDecimalDigits": 3
}
\end{verbatim}

\section*{6. Recomendaciones de uso}

\begin{itemize}
\item Verificar siempre los puertos COM antes de iniciar la aplicaci\'on.
\item Confirmar que la configuraci\'on serial coincide con la del hardware.
\item Ajustar correctamente \texttt{WeightDecimalDigits} y \texttt{TareDecimalDigits} seg\'un la balanza instalada.
\item Revisar los logs si no se reciben datos o si una pesada no se guarda.
\item No modificar directamente la base local salvo por necesidad t\'ecnica.
\end{itemize}

\section*{7. Logs y diagn\'ostico}

La aplicaci\'on genera logs en la carpeta:

\begin{itemize}
\item \texttt{logs}
\end{itemize}

Nombre del archivo:

\begin{itemize}
\item \texttt{gestion-de-fardos-YYYYMMDD.log}
\end{itemize}

All\'i se registran:

\begin{itemize}
\item inicio y cierre de la app;
\item apertura y cierre de puertos;
\item datos recibidos en crudo;
\item interpretaci\'on de balanza;
\item eventos del pulsador;
\item guardados;
\item ediciones;
\item exportaciones;
\item borrados hist\'oricos;
\item errores.
\end{itemize}

\section*{8. Resumen final}

La aplicaci\'on \textbf{Gesti\'on de Fardos} permite operar el proceso completo de captura:

\begin{itemize}
\item leer el peso de una balanza;
\item registrar la pesada al presionar el pulsador;
\item guardar el resultado en una base local;
\item corregir registros;
\item exportar datos a Excel;
\item diagnosticar la comunicaci\'on desde un modo t\'ecnico.
\end{itemize}

La configuraci\'on m\'as importante para esta versi\'on es la interpretaci\'on decimal del peso y de la tara, ya que el equipo informa los valores como enteros de 6 d\'igitos y la posici\'on de la coma debe definirse manualmente en \texttt{config.json}.

\end{document}
